\subsection{Tests and Builds}
\label{subsec:tests-builds}

The governance separates repository-specific rules from established language
tools. Scanning, severity, rule IDs, exclusions, exceptions, architecture
rules, and the aggregate report remain in the shared quality package; linting,
type checking, and formatting are delegated to the respective toolchains.
Table~\ref{tab:quality-language-tools} summarizes the coverage.

\begin{table}[H]
  \centering
  \scriptsize
  \renewcommand{\arraystretch}{1.08}
  \caption{Language-specific tools and boundaries in the quality gate}
  \label{tab:quality-language-tools}
  \begin{tabularx}{\textwidth}{@{}p{0.13\textwidth}p{0.28\textwidth}>{\raggedright\arraybackslash}X p{0.25\textwidth}@{}}
    \toprule
    \textbf{Area} & \textbf{Tools} & \textbf{Automated} &
      \textbf{Principal limitation} \\
    \midrule
    Python & Tokenizer, Python AST, Ruff & LOC, function and class size,
      McCabe complexity, nesting, parameters, lint, and format & Syntax errors
      can omit scope metrics in a focused size run. \\
    JS/TS & Dedicated LOC lexer, TypeScript AST, ESLint, TypeScript compiler,
      Prettier & LOC, classes and imports, function metrics, lint, types, and
      format & Architecture resolution covers only known relative and configured
      alias paths. \\
    Rust/Tauri & LOC lexer, Syn AST analyzer (WASI), Wasmtime, rustfmt,
      Clippy, Cargo & File and function size, complexity, nesting, parameters,
      format, lint, compiler check & Syn analyzes syntax without type resolution
      or macro expansion; adapter semantics remain subject to review. \\
    \bottomrule
  \end{tabularx}
  \smallskip
  {\footnotesize Source: Author's own compilation based on
  \textcite{kleiveist2026quality-implementation}.\par}
\end{table}

Ruff is declared as \texttt{>=0.14,<1.0}; version~0.16.4 was run locally. The
frontend lockfile pins ESLint~9.39.5, Prettier~3.9.6, TypeScript~5.9.3, and
typescript-eslint~8.67.0; the release workflows use Node~24. The Tauri
application manifest names Rust~1.77 and edition~2021 as its minimum. Separately,
the analyzer build environment pins rustc~1.97.1 and edition~2024. Its lockfile
pins Syn~2.0.119 and proc-macro2~1.0.107, while the tooling environment pins
Wasmtime~47.0.1. These versions improve snapshot traceability but remain time-
and environment-dependent.

The scanner recognizes \path{.py}, \path{.ts}, \path{.tsx}, \path{.js},
\path{.jsx}, \path{.mjs}, \path{.cjs}, and \path{.rs}. Physical lines are
distinguished from code LOC: blank lines and comment-only lines do not count,
while code followed by a comment counts once. Python docstrings and other
multiline strings count as code. For JavaScript/C-like syntax, a dedicated
lexer handles block comments; the Rust LOC pass also handles nested block
comments and raw strings. This lexer stage classifies code lines only. Rust
scope and control-flow metrics are parsed separately from a complete syntax
tree by the analyzer built with Syn~2.0.119. \path{.mjs} and \path{.cjs}
receive file LOC and the usual applicable frontend checks, but no
repository-specific scope metrics or architecture checks.

The checked-in analyzer targets \texttt{wasm32-wasip1} and is executed by
Wasmtime~47.0.1 from the tooling environment. The 535,787-byte artifact has
SHA-256
\path{7a27cb02a9392b62c487b4ce73a03524d7260b804c0c49eb4520ceb1a1cacfd8}.
Its provenance binds the exact rustc 1.97.1 compiler commit, target and
dependency versions, and a source digest over \path{build.py},
\path{Cargo.toml}, \path{Cargo.lock}, \path{rust-toolchain.toml}, and
\path{src/**/*.rs}.

The build removes inherited compiler, flag, target, and release-profile
overrides. A versioned remapping contract canonicalizes the source root, user
home, Cargo home and target, Rust sysroot, and every dependency root; broad
mappings precede specific ones because the last matching rustc entry wins.
Residual private physical build paths in the artifact cause failure. Before
each run, the host checks the provenance and artifact. The WASI instance has
bounded memory and fuel; malformed or unsupported syntax, invalid provenance,
a corrupt artifact, resource exhaustion, and inconsistent JSON metrics abort
the analysis. Generated profiles without a Rust toolchain can therefore run
the same Rust analysis. The analyzer replaces neither rustc nor Clippy: it
does not resolve types, expand macros, or check runtime semantics.

Excluded directories are \path{.cache}, \path{.dist}, \path{.generated},
\path{.git}, \path{.pytest_cache}, \path{.venv}, \path{__pycache__},
\path{coverage}, \path{dist}, \path{generated}, \path{node_modules},
\path{target}, and \path{vendor}; \path{Cargo.lock},
\path{package-lock.json}, and the configured paths \path{build/**},
\path{backend/build/**}, \path{frontend/dist/**}, \path{src-tauri/gen/**}, and
\path{src-tauri/target/**} are also excluded. The general directory name
\path{build} is intentionally not excluded globally, so handwritten tooling
paths with that name remain visible. Exclusions define the scope, for example
for generated or external artifacts; patterns that are too broad could hide
relevant code.

Exceptions are separate, temporary deviations for handwritten files. Required
fields are a known CQ or AR identifier, an exact relative path, a rationale of
at least ten characters, and an ISO expiry date; an optional symbol is matched
exactly. The file must exist, and duplicates and unknown IDs are rejected. The
expiry date is inclusive; only the following day produces \texttt{EX002}, and
a suppression remains visible in the report. The implementation does not,
however, check whether the scanner covers the file or whether an optional
symbol actually exists in the source. Further discrepancies are discussed in
Section~\ref{subsec:schwaechen-grenzen}.

In the backend, a Python AST analysis maps imports to four layers: API,
Application/Services, Domain, and Infrastructure. It detects configured
forbidden directions, direct Domain dependencies on seven framework roots,
and direct database imports in routers. Handlers above 50 code LOC also
produce \texttt{AR004}. In the frontend, boundaries between Application,
Features, API, UI, and Shared are checked, together with cross-feature imports
of internals; the target feature's public root \path{index.*} remains allowed.
Statically resolvable \texttt{import()} and \texttt{require()} calls are
included; dynamic and transitive dependencies that cannot be resolved
statically and unknown aliases remain outside the analysis. A router heuristic
cannot fully identify semantic business logic; thin Tauri commands are a
review-only rule.

The eight test suites cover tooling, schema, API, database, PostgreSQL,
frontend, E2E, and Tauri/Rust. \texttt{test --suite all} provides the complete
project-wide entry point; reports are optional
\parencite{kleiveist2026control,kleiveist2026tooling}.
Among other elements, the tooling suite tests the generator, profiles, and
configuration; the component suites delegate to Pytest, Vitest, Playwright, or
Cargo. The test areas therefore remain independently executable.

Selection remains profile-dependent: if a feature is absent, the suite reports
\texttt{SKIP}; if sources or prerequisites are missing for an active feature,
it reports \texttt{FAIL}. PostgreSQL without a test URL and E2E without
Playwright configuration are skipped. The runner can repair missing backend
environments and manage the development services for E2E.
This behavior deliberately distinguishes absent features from defective active
components. A skipped test is therefore not evidence of success.

The build paths remain separate: web produces Vite output and a ZIP; desktop
delegates the target and bundle to Tauri; and container builds frontend,
backend, or both images for cloud profiles. \texttt{container validate} checks
the Compose model without starting it
\parencite{kleiveist2026container,kleiveist2026release}.
The web path additionally creates a ZIP package. Desktop builds remain tied to
an active Tauri feature and the respective target platform.

GitHub Actions uses the same public entry points for the Core, profile
matrices, PostgreSQL, and desktop targets. CI supplements these with operating
systems, runtimes, and a temporary PostgreSQL service
\parencite{kleiveist2026ci}.
Local and external workflows therefore use comparable entry points, but do not
run under identical environmental conditions.
\beginnerinput{workspace/statement/500_tooling_workflow/550}
